\chapter{Paire de Planck, tension universelle et compliance gravitationnelle}
\label{chap:pareja-planck-respuesta}

\section{Paire dimensionnelle de Planck \((Q,\ell)\) : \(Q\ell=\hbar c\), \(\ell/Q=G/c^4\) et reconstruction univoque}

Soient \(t_0\) l'unité chronologique de la section CT108 et \(c\) la
section causale sélectionnée par le lecteur constitutif. Nous définissons

\[
t_P=\frac{54}{\pi}t_0,
\qquad
\ell=\ell_P=ct_P=\frac{54}{\pi}ct_0,
\]

et le quantum énergétique

\[
Q=E_P=\frac{\hbar}{t_P}=\frac{h}{108t_0}.
\]

La périodicité \(T_{108}=108t_0\) vérifie

\[
T_{108}=2\pi t_P,
\qquad
QT_{108}=h.
\]

De ces définitions se déduit la première identité de la paire :

\begin{equation}
\boxed{Q\ell=\hbar c.}
\label{eq:rm-planck-product}
\end{equation}

La réalisation gravitationnelle circulaire détermine

\[
G=\left(\frac{54}{\pi}\right)^2\frac{c^5t_0^2}{\hbar}.
\]

En substituant \(\ell=(54/\pi)ct_0\) et \(Q=\hbar/t_P\), on obtient la
deuxième identité :

\begin{equation}
\boxed{\frac{\ell}{Q}=\frac{G}{c^4}.}
\label{eq:rm-planck-quotient}
\end{equation}

Les équations \eqref{eq:rm-planck-product} et
\eqref{eq:rm-planck-quotient} sont algébriquement indépendantes comme
relations de produit et de quotient. Elles reconstruisent de façon univoque

\[
Q=\sqrt{\frac{\hbar c^5}{G}},
\qquad
\ell=\sqrt{\frac{G\hbar}{c^3}}.
\]

\begin{theorem}[Clôture de la paire de Planck]
\label{thm:rm-planck-pair}
Dans la réalisation circulaire CT108, la paire positive \((Q,\ell)\) est
l'unique solution de
\[
Q\ell=\hbar c,
\qquad
\ell/Q=G/c^4.
\]
Elle satisfait, en outre,
\[
GQ=c^4\ell,
\qquad
GQ^2=\hbar c^5,
\qquad
G\hbar=c^3\ell^2.
\]
\end{theorem}

\begin{proof}
La positivité permet de multiplier et de diviser les deux équations sans
ambiguïté de signe. La première égalité supplémentaire provient de
\(\ell/Q=G/c^4\) ; la deuxième résulte de sa multiplication par \(Q^2\), et la
troisième de la multiplication de \(\ell/Q=G/c^4\) par
\(Q\ell=\hbar c\). La reconstruction par racines positives démontre
l'unicité.
\end{proof}

\section{Réciprocité entre tension et compliance de Planck et conservation du produit d'action}

La paire définit deux réponses réciproques :

\begin{equation}
\mathscr T_P:=\frac{Q}{\ell}=\frac{c^4}{G},
\qquad
\mathscr C_P:=\frac{\ell}{Q}=\frac{G}{c^4}.
\label{eq:rm-tension-compliance}
\end{equation}

La première possède une dimension de force ; la deuxième, une dimension de
compliance gravitationnelle. Elles satisfont

\[
\mathscr T_P\mathscr C_P=1.
\]

L'équation d'Einstein s'écrit alors comme une loi constitutive :

\begin{equation}
\boxed{
G_{\mu\nu}+\Lambda g_{\mu\nu}
=8\pi\mathscr C_P T_{\mu\nu}.}
\label{eq:rm-einstein-compliance}
\end{equation}

La matière apporte une tension énergétique et la géométrie répond au moyen de
\(\mathscr C_P\). La constante \(G\) publie la capacité de courbure par
unité de tension dans la carte causale choisie.

Dans le régime newtonien, pour une densité énergétique
\(\varepsilon\),

\[
\nabla^2\!\left(\frac{\Phi}{c^2}\right)
=4\pi\mathscr C_P\varepsilon.
\]

L'action d'Einstein--Hilbert est normalisée par l'aire structurelle :

\begin{equation}
\frac{S_{\rm EH}}{\hbar}
=\frac{1}{16\pi\ell^2}
\int_{\mathcal M}(R-2\Lambda)\sqrt{-g}\,\mathrm d^4x.
\label{eq:rm-eh-area}
\end{equation}

La réponse scalaire et la réponse tensorielle sont des publications différentes
de la même compliance. La première utilise \(4\pi\mathscr C_P\) ; la deuxième,
\(8\pi\mathscr C_P\). En insérant CT108,

\[
2\pi\mathscr C_P=108\frac{ct_0}{Q},
\qquad
4\pi\mathscr C_P=216\frac{ct_0}{Q},
\qquad
8\pi\mathscr C_P=432\frac{ct_0}{Q}.
\]

La progression \(108\to216\to432\) se déduit des facteurs géométriques.
La coïncidence de \(432\) avec la multiplicité de la région transversale
de \(\pi\) constitue une compatibilité structurée ; sa promotion exige
une application d'incidence qui transporte la multiplicité généalogique vers le tenseur
d'Einstein.

\section{Quantum informationnel ternaire \(\log 3\) et capacité de la cellule qutrit}

La relation thermodynamique de la section CT108 est

\begin{equation}
\boxed{Q=k_B\Theta_{\rm clk}\log3.}
\label{eq:rm-planck-trit}
\end{equation}

Dans cette équation, \(\log3\) est l'entropie d'une alternative ternaire
équiprobable ; \(\Theta_{\rm clk}\) fixe la carte thermique et \(Q\) publie
l'énergie du cycle. La combinaison avec le théorème
\ref{thm:rm-planck-pair} produit

\begin{equation}
\boxed{
Gk_B\Theta_{\rm clk}\log3=c^4\ell.}
\label{eq:rm-gravity-information}
\end{equation}

La gravité transforme le quantum informationnel de l'horloge en une longueur de
réponse géométrique. De manière équivalente,

\[
\ell=\mathscr C_P k_B\Theta_{\rm clk}\log3.
\]

Cette égalité sépare le contenu adimensionnel \(\log3\), le cadre
thermométrique \(k_B\Theta_{\rm clk}\) et la compliance géométrique.

\section{Covariance des lecteurs bidirectionnels sous inversion de feuillet et conjugaison temporelle}

Soient \(\hbar_{\rm pre}\) et \(\hbar_{\rm ret}\) les deux sections orientées
de la droite d'action, reliées par un quotient positif
\(R_{\rm act}\) :

\[
\frac{\hbar_{\rm pre}}{\hbar_{\rm ret}}=R_{\rm act}.
\]

L'invariance de \(\ell\) impose

\[
\frac{G_{\rm pre}}{G_{\rm ret}}=R_{\rm act}^{-1}.
\]

Par conséquent,

\begin{equation}
G_{\sigma}\hbar_{\sigma}=c^3\ell^2
\qquad(\sigma\in\{\mathrm{pre},\mathrm{ret}\}).
\label{eq:rm-two-way-area}
\end{equation}

L'action conserve l'orientation du chemin ; le produit par la réponse
gravitationnelle conserve l'aire. Cette covariance fait de \(\ell^2\)
l'invariant géométrique de l'inversion bidirectionnelle.

\section{Vide constitutif et clôture causale}

L'opérateur de vide fournit

\[
c_{\rm vac}^{-2}=\mu\varepsilon.
\]

La paire de Planck fournit une deuxième section causale :

\[
c_{\rm int}=\frac{Q\ell}{\hbar}.
\]

Nous définissons le scalaire de raccordement

\begin{equation}
\Xi_c=\frac{Q\ell}{\hbar}\sqrt{\mu\varepsilon}
=\frac{c_{\rm int}}{c_{\rm vac}}.
\label{eq:rm-causal-seam}
\end{equation}

L'égalité \(\Xi_c=1\) est nécessaire et suffisante pour que les deux lecteurs
publient le même cône causal. Sous ce raccordement,

\begin{equation}
\boxed{
G=\left(\frac{54}{\pi}\right)^2
\frac{t_0^2}{\hbar(\mu\varepsilon)^{5/2}}.}
\label{eq:rm-g-vacuum-product}
\end{equation}

\begin{proof}
La relation \(c=(\mu\varepsilon)^{-1/2}\) donne
\(c^5=(\mu\varepsilon)^{-5/2}\). Sa substitution dans le sélecteur circulaire de
\(G\) produit \eqref{eq:rm-g-vacuum-product}.
\end{proof}

Le produit \(\mu\varepsilon\) gouverne le cône causal et la réponse
gravitationnelle ; le quotient \(\mu/\varepsilon=Z^2\) gouverne l'impédance et
l'orientation. Les deux contenus proviennent du même opérateur bicouche et conservent
des fonctions distinctes.

\section{Factorisation machienne de la réponse inertielle locale au moyen de la trace de la clôture globale : \(\mathcal I_v=\overline{\mathcal I}_v\circ b\)}

Une énergie locale \(E\) détermine la coordonnée adimensionnelle

\[
\xi_E(E)=\frac{E}{Q}.
\]

Le couplage gravitationnel de deux énergies est

\begin{equation}
\alpha_G(E_1,E_2)=\frac{E_1E_2}{Q^2}.
\label{eq:rm-mach-coupling}
\end{equation}

La référence \(Q\) appartient à l'horloge globale de clôture. L'inertie locale
s'évalue donc par rapport à une échelle produite par la généalogie complète.
Cette dépendance ne postule pas une action instantanée de l'univers sur le
corps : elle établit une factorisation opératoire de la réponse locale par la
classe globale de clôture.

\begin{proofstatusbox}
Les identités de la paire, la tension, la compliance et la covariance
\(G\hbar\) sont exactes dans la réalisation circulaire CT108 déclarée.
L'identification \(c_{\rm vac}=c_{\rm int}\) est une interface physique typée ;
son contenu se concentre dans l'application de droites causales dont l'évaluation est
\(\Xi_c=1\).
\end{proofstatusbox}

\begin{falsifierbox}
Réfutent la clôture de ce chapitre : une autre paire positive qui
satisfasse simultanément les deux équations directrices ; la perte de
\(G\hbar=c^3\ell^2\) sous l'inversion d'action ; une section de vide
qui ne satisfasse pas \(c^{-2}=\mu\varepsilon\) ; ou un transport de cadre pour lequel
\(\Xi_c\ne1\).
\end{falsifierbox}
